\documentclass[11pt,fleqn]{report}
\usepackage{graphicx}
\usepackage{amssymb}
\usepackage{longtable}
\usepackage[T1]{fontenc}

\usepackage{cml}
\lstnewenvironment{codeXML}
    {\lstset{basicstyle=\small\ttfamily,
             language=XML,
             xleftmargin=0.0in,
             xrightmargin=0.0in,
             backgroundcolor=\color{Orchid!60}, % 60% orchid, 40% white
             escapeinside={\%*}{*)}, %  .tex file characters inside %*...*) are
                                     %  treated as Tex commands.
            }
    }
    { }

\newcommand{\ModelName}{XML Helper}
\newcommand{\ModelAuthor}{Daniel Ghan \\ Gary Turner}
\newcommand{\ModelAffiliation}{Odyssey Space Research}
\newcommand{\ModelDate}{July 2022}
\pagestyle{plain}
\usepackage[parfill]{parskip}
\usepackage[margin=1.5cm]{caption}
\usepackage{subfig}

\begin{document}
\titlePage
\clearpage % Reset page number
\pagenumbering{roman}
\tableofcontents % Print table of contents
\vfill
{\Large \textbf {Revision History}}
{\large
\begin{center}
\begin{tabular}{|l||l|l|l|}
\hline
 \textbf{Version} & \textbf{Date} & \textbf{Author} & \textbf{Purpose} \\ \hline
 \hline
1 & July 2022 & Gary Turner & Initial Version \\ \hline
\end{tabular}
\end{center}
}

\chapter{Introduction} % Make a "chapter" heading
\pagenumbering{arabic} % Start text with arabic 1
\par The XML-helper model provides some utility functions to facilitate
processing of XML data files. This model was written to support the
fault-injection / fault-management CML model which parses XML files to
specify the type and nature of faults to be injected. It can be used in
other applications.

\chapter {Requirements}
The model shall provide commonly used functions to support processing of XML
data. It is anticipated that this model is incomplete. The most urgent needs
identified are:
\begin{enumerate}
  \item  Provide the value of a specified node property for a specified node.
  In the example below, we may want the value of the property ID or purpose for
  this node.
  \begin{codeC}
    <ABC name="myname" purpose="none" ID="5"></ABC>
  \end{codeC}
  \item Locate an XML node by node-name (or node-type) within the node-tree.
  Three levels of specificity are required:
  \begin{enumerate}
    \item some specified level in the XML tree; this node would start at the
    next sibling of a specified node. This option subdivides a tree-level,
    looking only for nodes that are at the specified location, or downstream
    siblings of it. It cannot return "upstream" nodes from the specified level.
    \item the level one below some specified level in the XML tree; this node
    would be a child of a specified node. This option will search an entire
    tree-level until it finds a matching node. For a tree-level containing
    multiple nodes of the same type, this option will always return the
    first encountered, it cannot find additional downstream nodes of the
    same type.
    \item any level below some specified level in the XML tree; this node
    would be a progeny of unspecified generation of a specified node. This
    option will iteratively search through the tree, descending through
    generations first, and then across generations.
  \end{enumerate}
  Notes:
  \begin{itemize}
    \item A node-name in this context is the type of node, not one of the
    node properties. For example, in the XML line above, the node-name would
    be `ABC'; the field identified as `name' is a property of this node.
  \end{itemize}
\end{enumerate}



\chapter {Model Specification}
The model is written to provide a set of static methods; it is not intended
that this model be instantiated at any time.
The model provides the following four methods to address the four algorithms
specified in the requirements section.
\begin{codeC}
  static xmlNodePtr xml_find(        xmlNodePtr node, const char* name);
  static xmlNodePtr xml_find_child(  xmlNodePtr node, const char* name);
  static xmlNodePtr xml_find_progeny(xmlNodePtr node, const char* name);

  static const char* xml_find_value( xmlNodePtr node,
                                     const char* name,
                                     bool allow_case=false);
\end{codeC}
It also provides additional utilities for convenience, such as verification
that a node has the correct name / type in \textit{xml\_name\_match}.


\chapter {User's Guide}
The model methods are called as needed, for example:
\begin{codeC}
const char * value = XmlHelper::xml_find_value( node_A, "ID");
\end{codeC}

The five methods listed in the previous chapter are discussed below.

\begin{itemize}
\item \textbf{xml\_find} looks for nodes with the specified name; nodes must
be ``downstream'' sibling nodes of the specified node.
It returns a pointer to the node, or \textit{nullptr} if none is found.
\item \textbf{xml\_find\_child} looks for nodes with the specified name;
nodes must be immediate children of the specified node.
It returns a pointer to the node, or \textit{nullptr} if none is found.
\item \textbf{xml\_find\_progeny} looks for nodes with the specified name;
nodes must be subordinate to the specified node.
It returns a pointer to the node, or \textit{nullptr} if none is found.

While \textit{xml\_find\_child} will search only one generation down,
\textit{xml\_find\_progeny} will search the entire tree below the specified
node.

\item \textbf{xml\_find\_value} looks for a field with the specified name
within the specified node. It returns a \textit{char*} pointer to the array
containing the string found at the named field, or \textit{nullptr} if no
field is found with a matching name.

The optional flag supports flexibility where the field-name may have a first
character with inconsistent case. For example, in the verification tests, we
have \textit{Node2} with a field-name \textit{Purpose}, while the other nodes
have a field-name \textit{purpose}. Searching \textit{Node2} for
\textit{purpose} with:
\begin{codeC}
const char* purpose =XmlHelper::xml_find_value( node2,"purpose");
\end{codeC}
will default to requiring a perfect string comparison, and no match will be
found. However, using:
\begin{codeC}
const char* purpose =XmlHelper::xml_find_value( node2,"purpose",true);
\end{codeC}
will support a positive string match when comparing \textit{purpose} and
\textit{Purpose}. Note that this case insensitivity option:
\begin{itemize}
\item is limited to field-names, it does not apply to node-names;
\item is limited to the first character of a field-name;
\item must be activated by providing \textit{true} as the optional third
argument -- if only two arguments are provided, the field-names must match
the specified name completely.
\end{itemize}

\item \textbf{xml\_name\_match} is used by \textit{xml\_find},
\textit{xml\_find\_child}, and \textit{xml\_find\_progeny} to compare the
node-names against the specified name. For this name comparison, a case
match is required on all characters.
\end{itemize}


\chapter{Verification}

\section{Code Coverage}
\begin{figure}[h]
  \centering
  \includegraphics[scale=0.6]{graphics/code_coverage.png}
  \caption{Code Coverage}
  \label{fig:coverage}
\end{figure}
\newpage

\section {Coding Standards}
There is a coding-standards violation in the model, in the use or
\textit{reinterpret\_cast}. This model will need a waiver to allow the use this
dangerous function. The specific use-case is in reinterpreting a pointer to an
\textit{unsigned char} as a pointer to a \textit{char}.

The rationale for this need follows:
\begin{itemize}
  \item The *.xml files provide information about the nodes as ASCII strings.
  \item The *.xml files are parsed and processed by the XML library, which
  stores the content it reads as \textit{unsigned char} arrays.
  \item Other aspects of the simulation, including the SWIG interface and other
  user-settable configurations, typically use \textit{char} arrays to store the
  expected ASCII strings.
  \item We need to be able to compare the expected strings with those stored by
  the XML reader, so it is necessary to reinterpret one as the other.
  \item Compilers will not implicitly cast between \textit{unsigned char} and
  \textit{char} because one half of the 255 possible values of the 8-bit
  variables represent different interpretations.
  \item The g++ compiler interprets the values 0-127 of a \textit{char} variable
  equivalently to the values 0-127 of an \textit{unsigned char} variable.
  (the values 128-255 may be treated differently).
  NOTE THAT THIS IS NOT GUARANTEED BY THE C++ STANDARD.
  \item The ASCII text map uses values 0-127, so ASCII text may be equivalently
  stored as \textit{unsigned char} or \textit{char} with our compiler.
  \item This model is limited to comparing ASCII text, as read by the XML
  reader.
  \item In this model, the reinterpret cast is limited to converting from
  \textit{unsigned char} to \textit{char}; it is isolated to a method
  that requires an input argument of type \textit{unsigned char*} and outputs a
  \textit{char*}.
  \item The verification process includes comparison between values provided as
  \textit{char} with values stored by the XML-reader as \textit{unsigned char}.
  If the compiler were to change at some point, and if that change resulted in
  breaking the current equivalence between these types for ASCII characters,
  then the verification tests would fail. There is no reasonable path that
  might lead to a situation in which the verification tests pass and the model
  creates a false positive match between a user-specified
  \textit{char} string and an XML-read \textit{unsigned char} string.
\end{itemize}

\section {Verification Tests}
The model provides a single test case, at \textit{SIM\_verif/RUN\_verif}.
This test case exercises all of the functionality of the model, using an
XML structure defined in sample.xml with nodes as illustrated below.
The structure includes four nodes, colored orange, that are all of the
same type (\textit{TargetNode}); these will be the target of the search
algorithms.
\begin{figure}[h]
  \centering
  \includegraphics[scale=0.6]{graphics/verif_tree.png}
  \caption{Node Tree from sample.xml}
\end{figure}

The simulation utilizes a test class \textit{XmlHelperTester}, provided at
\textit{SIM\_verif/include/test.hh}.
Most of the verification is conducted within \textit{XmlHelperTester::test()}.
The execution of this test relies on internal testing; with each test if
the results are unexpected, the simulation will terminate prematurely with
a failed flag. There is no regression data to compare; if any tests fail,
the resulting premature termination will indicate that failure.

The verification includes several tests:
\begin{enumerate}
  \item Testing name-match (\textit{xml\_name\_match})
  \begin{enumerate}
    \item Tests the root-node to confirm that its name matches with expectation.
  \end{enumerate}
  \begin{figure}[h]
    \centering
    \includegraphics{graphics/snapshot1.png}
  \end{figure}

  \item Testing property lookup (\textit{xml\_find\_value}):
  \begin{enumerate}
    \item parses through the level 1 nodes and prints their \textit{purpose}
    values.  These tests do not have individual fail-checks, but the test
    will fail if it does not detect 4 nodes at this level.
    \item The verification sequence identifies some useful nodes for later,
    and uses the \textit{xml\_find\_value} method to confirm that the correct
    nodes have been identified, printing the output to the screen.
    This serves no purpose for verification, it is informational as a
    troubleshooting tool.
    If the wrong nodes are identified, later tests will fail.
    \item Note that we deliberately introduced a case error in the
    field name \textit{purpose} for \textit{Node2}, labeling it \textit{Purpose}
    instead. This was done to test the optional case-insensitivity setting
    on the first character of a field name. Consequently, everywhere that
    \textit{xml\_find\_value} is used with \textit{Node2}, we must use the
    optional flag to allow the field-name \textit{Purpose} to match with the
    specified name, \textit{purpose}.
  \end{enumerate}
  \begin{figure}[h]
    \centering
    \includegraphics{graphics/snapshot2.png}
  \end{figure}

  \item Testing the same-level node finder (\textit{xml\_find})
  \begin{enumerate}
    \item Starting with node-1, the search terminates at node-2,
    the first downstream sibling TargetNode.
    \item  Starting with node-3, the search terminates at node-4, the
    first downstream sibling TargetNode; even though node-2 is also a
    sibling, it is upstream of the starting-point (node-3) and so is
    not considered.
    \item  Starting with node-5, the search terminates with no result;
    even though node-8 is on the same level, and "downstream", it is not
    a sibling of node-5.
  \end{enumerate}
  \begin{figure}[h]
    \centering
    \includegraphics{graphics/snapshot3.png}
  \end{figure}

  \item Testing the child-level node-finder (\textit{xml\_find\_child})
  \begin{enumerate}
    \item Starting with node-0, the search terminates at node-2,
    the highest-upstream of the child nodes with type \textit{TargetNode}.
    \item Starting with node-2, the search terminates with no result;
    neither node-5 nor node-6 are \textit{TargetNode} nodes, and this
    search will only include immediate child nodes, it will not descend
    further.
    \item Starting with node-3, the search terminates at node-8,
    the highest-upstream of the child nodes with type \textit{TargetNode}.
  \end{enumerate}
  \begin{figure}[h]
    \centering
    \includegraphics{graphics/snapshot4.png}
  \end{figure}

  \item Testing the tree-search node-finder (\textit{xml\_find\_progeny})
  \begin{enumerate}
    \item Starting with node-0, the search terminates at node-2 after
    considering node-1.
    \item Starting with node-2, the search terminates at node-11.
  \end{enumerate}
  \begin{figure}[h]
    \centering
    \includegraphics{graphics/snapshot5.png}
  \end{figure}
\end{enumerate}

\end{document}
